\section{Limitations}
\label{sec:limitations}

\paragraph{One laptop.}
Every measurement comes from one Apple M2 Max with 32~GB of memory, and the machine bounds the study. The weight-level methods, exact scoring, probes and LoRA, need the weights in the local backend, so they stop at 8B parameters, and LoRA, limited by training time, stops at 4B with 40 training episodes. Our scale slopes therefore describe the range from 135M to 8B and cannot say whether the curves bend, flatten or jump above it. The H3 pairs are 12 to 15 times apart in parameters, not the hundredfold spans that larger ladders allow, and the reasoning ablation stops at 256 tokens, a budget that Qwen3's thinking mode may use up before it reaches an answer. Section~\ref{sec:future} lists the rungs and conditions that did not fit. We think the restriction is a trade worth making, since anyone with similar hardware can rerun the whole study, but it is a restriction.

\paragraph{Synthetic tasks.}
The seven tasks are small: at most ten actions, at most 100 decisions per episode, and states that fit in a few lines of text. We chose them because each has an exact oracle and a distinct decision structure, not because they resemble deployed agentic work, which involves long contexts, tool outputs and open-ended actions. Slopes measured here need not transfer to such settings, and a model's rank on \deadeye{} is evidence about the decision structures we test, not about agentic ability in general.

\paragraph{Oracles.}
Every oracle is optimal for its task, but several properties of the oracles limit what agreement and cloning can show. They break ties by fixed rules, so oracle agreement understates the share of optimal decisions wherever several moves are equally good. The blackjack oracle is optimal only for a reduced game without doubling, splitting or surrender, so our blackjack is not the casino game. The two bandit oracles are clairvoyant: a normalized score of 1 is out of reach there, and the labels used to train probes and LoRA come from a teacher whose choices cannot be inferred from the observation. A learnable teacher such as UCB1 would make behavior cloning on the bandits meaningful \citep{laskin2023incontext,nie2025evolve}; v1 does not use one.

\paragraph{Scale within families only.}
We draw scale conclusions only within families, and even there parameter count is confounded with the number of training tokens per parameter \citep{hoffmann2022chinchilla,kumar2025precision}. The ladders are short, three or four rungs each, so a slope summarizes a trend more than it establishes a functional form, and its bootstrap interval reflects the noise of episodes, not the idiosyncrasies of individual checkpoints. Every rung runs with 16-bit weights, so scale is not mixed with quantization, but total parameter counts overstate the capacity of the smallest models, in which embeddings take a large share. Comparisons across families are descriptive.

\paragraph{Prompt sensitivity.}
Each condition uses one frozen prompt. Changes of format that preserve meaning can move accuracy by tens of points \citep{sclar2024formatting}, and our format ablations, letters against names, demonstrations, temperature and length normalization, cover only a few such changes on one ladder. We do not measure sensitivity to wording, and a different prompt might move individual cells by more than the effects we test.

\paragraph{Quantization.}
The quantization ablation measures two GGUF formats, Q8\_0 and Q4\_K\_M, as llama.cpp runs them. It says nothing directly about the NF4 and LLM.int8() weights of bitsandbytes \citep{dettmers2023qlora,dettmers2022llmint8}, which much of the work on quantized language models uses, or about GPTQ and AWQ \citep{frantar2023gptq,lin2024awq}. The 16-bit reference runs in PyTorch and the quantized models in llama.cpp, so each difference also contains whatever separates the two runtimes, such as kernel numerics and the rendering of the chat template; serving a 16-bit GGUF file of the same checkpoint through the same server would separate the two. Only generation and first-token letter scoring can be compared, on one ladder at three sizes.

\paragraph{Scoring through a server.}
The quantized models are served through an OpenAI-compatible interface and scored by the log-probability of the first token of each label, taken from the top 20 alternatives. This is exact only when the labels start with distinct tokens, which is why the block uses letters, and letter labels bring option-identifier priors with them \citep{zheng2024selectors}. These rows are reported on their own and never enter the scale curves.

\paragraph{Contamination.}
Instances are generated from seeds and never appear on the web, but the rules and strategies of tic-tac-toe and blackjack are widely published, so part of a model's performance there may be recall of strategy rather than reasoning. The contextual bandit and the sign-flipped loan task are the cases in which prior knowledge cannot help or actively misleads.

\paragraph{Adaptation.}
Our fine-tuning is behavior cloning with LoRA, at one rank and one learning rate for every rung, on 40 training episodes in the main configuration and only up to 4B. The 7B and 8B rungs are never fine-tuned, so the LoRA half of H2b and all of H3a rest on the smaller rungs. Reinforcement-learning fine-tuning, which generalizes differently \citep{chu2025sft} and repairs some exploration failures \citep{schmied2026greedy}, is not part of v1, so our conclusions about small fine-tuned models concern cloning only. Out-of-distribution performance is tested on one task with two kinds of shift, which says nothing about shifts in task structure.

\paragraph{Statistics.}
One hundred episodes per cell detect a difference of 0.10 with power 0.8 only when the standard deviation of the paired differences is at most about 0.36 (Section~\ref{sec:setup:stats}), and the 50 episodes of the ablations and the free-reply configuration detect only about 0.14 at that spread. On tasks whose anchors lie close together, blackjack above all, power is lower even with twenty hands per episode, and a null result there is weak evidence of no effect; we report the minimum detectable difference wherever it exceeds the effect of interest. The anchors themselves are means over the evaluation seeds and we treat them as fixed, so their sampling error is not propagated into the confidence intervals of \norm{}; the paired tests do not depend on the anchors.

\paragraph{Decision models.}
The comparison with decision models is exploratory and partial. It includes only the systems that run on the laptop, Kev at 0.8B and 4B, Laya and SemIf. Jev, whose weights, architecture and training data are closed, the hosted Clef, Clef-flash and Perplexity decider, and the larger Kev checkpoints are left out, so the comparison cannot say how the systems that defined the category fare on our tasks. The open-weight decision models were trained on classification panels, not on our tasks, so a gap in either direction may reflect training data as much as method.

\paragraph{Latency on one runtime.}
Latency depends on the hardware, the batch size and the implementation. Ours comes from PyTorch's Metal backend, which is slower than runtimes built for Apple silicon such as MLX and llama.cpp. A model or method that is twice as fast here need not be twice as fast on a GPU or in another runtime, so our cost comparisons hold for this laptop and this runtime, and token counts are the better guide elsewhere. The quantized models and the decision models run in other runtimes on the same machine, so their latencies are not directly comparable with ours.
